% Generated draft: requires mathematical and source review.
\documentclass[11pt]{article}
\usepackage[margin=25mm]{geometry}
\usepackage{fontspec}
\setmainfont{Latin Modern Roman}
\usepackage{amsmath,amssymb,mathtools,unicode-math}
\setmathfont{Latin Modern Math}
\usepackage{enumitem,needspace,xurl,hyperref}
\setlength{\parindent}{0pt}
\setlength{\parskip}{5pt}
\setlength{\emergencystretch}{3em}
\newcommand{\N}{\mathbb{N}}
\newcommand{\Z}{\mathbb{Z}}
\newcommand{\Q}{\mathbb{Q}}
\newcommand{\R}{\mathbb{R}}
\newcommand{\C}{\mathbb{C}}
\newcommand{\F}{\mathbb{F}}
\newcommand{\A}{\mathbb{A}}
\newcommand{\PP}{\mathbb{P}}
\newcommand{\E}{\mathbb{E}}
\newcommand{\eps}{\varepsilon}
\newcommand{\dd}{\,\mathrm{d}}
\begin{document}
\section*{UnsolvedMath ID: 701453}
\textbf{Source problem code:} catalogue entry
\subsection{Definitions and mathematical statement}
Let \(F\) be a free group, with no restriction on its rank specified in the supplied record. A function \(q:F\to\R\) is a quasimorphism if its defect
\[
D(q):=\sup_{g,h\in F}|q(gh)-q(g)-q(h)|
\]
is finite. Thus \(D(q)\) is the least nonnegative real number bounding all these expressions. Under the standard interpretation of homogeneity, \(q\) is homogeneous if
\[
q(g^n)=nq(g)\qquad(g\in F,\ n\in\Z).
\]
The source omits the domain and quantifier for \(n\); the integer-power convention is made explicit here and requires confirmation against the original problem.

Two functions \(q_1,q_2:F\to\R\) are at bounded distance if
\[
\sup_{g\in F}|q_1(g)-q_2(g)|<\infty.
\]
The supplied statement asserts that, for every homogeneous quasimorphism \(\phi:F\to\R\), every quasimorphism \(\psi:F\to\R\) at bounded distance from \(\phi\) satisfies
\[
D(\psi)\geq \tfrac12D(\phi).
\]
The question is whether, for every free group \(F\) and every homogeneous quasimorphism \(\phi:F\to\R\), there exists a quasimorphism \(\psi:F\to\R\), not required to be homogeneous, such that
\[
\sup_{g\in F}|\phi(g)-\psi(g)|<\infty
\qquad\text{and}\qquad
D(\psi)=\tfrac12D(\phi).
\]
The bounded-distance condition in this concluding question is inherited from the preceding sentence of the source. The unrestricted quantification over free groups reflects the absence of a stated rank hypothesis.
\subsection{Short English statement}
Given a homogeneous quasimorphism on a free group, can one change it by a uniformly bounded function so that its defect becomes exactly half the original defect? The resulting quasimorphism need not be homogeneous. The integer-power interpretation of homogeneity and the absence of a rank restriction require source confirmation.
\subsection{Sources}
\begin{itemize}
\item Supplied local UnsolvedMath registry record: permanent ID \(701453\), exact source problem code \texttt{catalogue entry}, title \emph{Kourovka Notebook Problem 18.24}; upstream revision \texttt{d371ee917ac9647fbb06cd6fc529deec6f3dda82}.
\item Attributed in the supplied record to M. Burger and D. Calegari, proposed in 2014.
\item Supplied bibliographic attribution: \emph{The Kourovka Notebook (Unsolved Problems in Group Theory)}, 18th issue (2014), original problem 18.24; also identified with the 21st-issue edition (2026), arXiv:1401.0300v48. The collection description names E. I. Khukhro and V. D. Mazurov as editors. These references are reproduced from the supplied record and have not been independently consulted.
\end{itemize}
\end{document}
